%! TEX root = master.tex
\chapter{2}{Numbers and Parity}

Before we can prove anything about numbers, we have to agree on what they
are. Some of what follows comes from rather simple ideas, and some of it we
simply have to accept: there is no substitute for it, and without it we can
go no further. What we accept without proof is an \emph{axiom}. What we
merely agree to call something is a \emph{definition}.

\section{The number systems}\label{sec:number-systems}

We take for granted the natural numbers $\N = \set{1, 2, 3, \dots}$, the
integers $\Z = \set{\dots, -2, -1, 0, 1, 2, \dots}$ with their usual
arithmetic and order, and the real numbers $\R$, the points of the number
line.

\begin{axiom}[Well-ordering]\label{ax:well-ordering}
  Every nonempty set of natural numbers has a least element.
\end{axiom}

\begin{definition}[Rational and irrational numbers]\label{def:rational}
  A real number $q$ is \emph{rational}, written $q \in \Q$, if
  \[
    q = \frac ab \quad\text{for some } a, b \in \Z \text{ with } b \neq 0.
  \]
  A real number that is not rational is \emph{irrational}. The irrationals
  form the set $\R \setminus \Q$.
\end{definition}

So $\N \subset \Z \subset \Q \subset \R$, since $n = n/1$. Examples of
rationals are $-\frac{3}{13}$, $-1$, $0$, $2$ and $\frac{13}{2}$. The
condition $b \neq 0$ matters, because $\frac a0$ is not a number. It is not
obvious that irrational numbers exist at all; \cref{sec:sqrt2} shows that
$\sqrt2$ is one.

\begin{definition}[Complex numbers]\label{def:complex}
  Let $i$ be a number with $i^2 = -1$. A \emph{complex number} is one of the
  form
  \[
    z = a + ib, \qquad a, b \in \R,
  \]
  and the set of them is $\C$. Taking $b = 0$ shows that $\R \subset \C$.
\end{definition}

\begin{remark}
  Writing $i = \sqrt{-1}$ is common but risky. The rule
  $\sqrt{xy} = \sqrt x \sqrt y$, valid for $x, y \geq 0$, would give
  $-1 = i \cdot i = \sqrt{-1}\sqrt{-1} = \sqrt{1} = 1$. So $i^2 = -1$ is the
  safe definition.

  It is unfortunate that $i$ is called \emph{imaginary}. In truth all
  numbers are imaginary, in the sense that we made them up, and $i$ is no
  less real than $\sqrt2$ or $-1$.
\end{remark}

\section{Even and odd}\label{sec:parity}

\begin{definition}[Even and odd]\label{def:even-odd}
  An integer $a$ is \emph{even} if $a = 2b$ for some $b \in \Z$, and
  \emph{odd} if $a = 2b + 1$ for some $b \in \Z$.
\end{definition}

%! TEX root = ../master.tex
\begin{theorem}[Parity]\label{thm:parity}
  Every integer is either even or odd, and no integer is both.
\end{theorem}

\begin{proof}
  \emph{At least one.} Let $n \in \Z$ and let $q$ be the greatest integer with
  $2q \leq n$. Then $2(q + 1) > n$, so $r = n - 2q$ satisfies $0 \leq r < 2$.
  If $r = 0$, then $n = 2q$ is even. If $r = 1$, then $n = 2q + 1$ is odd.

  \emph{Not both.} If $n = 2a = 2b + 1$ with $a, b \in \Z$, then
  $2(a - b) = 1$. But $2(a - b)$ is $0$ if $a = b$, and at least $2$ in
  absolute value otherwise. So it is never $1$.
\end{proof}


Because of \cref{thm:parity}, ``not even'' and ``odd'' mean the same thing
for integers. Every integer is exactly one of the two: even \emph{exclusive
or} odd.

%! TEX root = ../master.tex
\begin{theorem}[Arithmetic of parity]\label{thm:parity-arithmetic}
  Let $m, n \in \Z$.
  \begin{enumerate}[label=(\alph*)]
    \item If $m$ and $n$ have the same parity, then $m + n$ is even. If they
      have different parities, then $m + n$ is odd.
    \item If $m$ is even, then $mn$ is even. If $m$ and $n$ are both odd,
      then $mn$ is odd.
  \end{enumerate}
\end{theorem}

\begin{proof}
  Write the even numbers as $2a$ and the odd ones as $2a + 1$.
  \begin{enumerate}[label=(\alph*)]
    \item $2a + 2b = 2(a + b)$ and $(2a + 1) + (2b + 1) = 2(a + b + 1)$ are
      even, while $2a + (2b + 1) = 2(a + b) + 1$ is odd.
    \item $(2a)n = 2(an)$ is even, and
      $(2a + 1)(2b + 1) = 2(2ab + a + b) + 1$ is odd. \qedhere
  \end{enumerate}
\end{proof}

%! TEX root = ../master.tex
\begin{theorem}[Parity of a square]\label{thm:square-parity}
  For $n \in \Z$, $n^2$ is even if and only if $n$ is even.
\end{theorem}

\begin{proof}
  If $n$ is even, then $n^2 = n \cdot n$ is even by
  \cref{thm:parity-arithmetic}. For the converse we prove the contrapositive
  (\cref{thm:contrapositive}): if $n$ is not even, then by \cref{thm:parity}
  it is odd, so $n^2$ is odd by \cref{thm:parity-arithmetic}, and hence not
  even, again by \cref{thm:parity}.
\end{proof}


\begin{exercise}\label{ex:consecutive}
  Prove that the product of two consecutive integers is even.
\end{exercise}

\begin{proof}
  Let the integers be $n$ and $n + 1$. By \cref{thm:parity}, $n$ is even or
  odd. If $n$ is even, then $n(n + 1)$ is even by
  \cref{thm:parity-arithmetic}(b). If $n$ is odd, then $n + 1$ is even by
  \cref{thm:parity-arithmetic}(a), and again the product is even.
\end{proof}

\begin{exercise}\label{ex:n2-n-1}
  Prove that $n^2 + n + 1$ is odd for every integer $n$.
\end{exercise}

\begin{proof}
  $n^2 + n = n(n + 1)$ is even by \cref{ex:consecutive}, so
  $n^2 + n + 1 = 2k + 1$ for some $k \in \Z$.
\end{proof}

\begin{remark}
  As a consequence, $x^2 + x + 1 = 0$ has no integer roots. Every integer
  input gives an odd output, and $0$ is even.
\end{remark}
